% ===== Document Preamble =====
\documentclass[11pt]{article}
\usepackage[margin=0.5in]{geometry}
\usepackage{amsmath,amssymb}
\usepackage{xcolor}
\usepackage{enumitem}
\usepackage{needspace}
\usepackage{comment}

% ===== Colors =====
\definecolor{AnsRed}{RGB}{170,0,0}
\definecolor{WrongBlue}{RGB}{0,0,160}

% ===== Choice list styling (tight, consistent) =====
\setlist[enumerate]{%
  leftmargin=1.5em,
  itemsep=0.25em,
  topsep=0.25em,
  label=\textbf{\Alph*.}
}

% ===== Style Wrappers =====
\newcommand{\QuestionStyle}[1]{{\large\bfseries\boldmath #1}}
\newcommand{\AnswerStyle}[1]{{\large\bfseries\boldmath #1}}
\newcommand{\ReasonStyle}[1]{{\normalsize #1}}

% ===== Show / Hide Reasoning Toggle =====
\newif\ifshowreason
\showreasontrue      % Instructor version
% \showreasonfalse   % Student version

% ===== Choice Item Macros (Answer + Reasoning) =====
% Usage:
%   \Correct{Answer}{Reasoning}
%   \Wrong{Answer}{Reasoning}

\newcommand{\Correct}[2]{%
  \item \AnswerStyle{\textcolor{AnsRed}{#1}}%
  \ifshowreason\\[-0.2em]
    \ReasonStyle{\textcolor{AnsRed}{#2}}%
  \fi
}

\newcommand{\Wrong}[2]{%
  \item \AnswerStyle{\textcolor{WrongBlue}{#1}}%
  \ifshowreason\\[-0.2em]
    \ReasonStyle{\textcolor{WrongBlue}{#2}}%
  \fi
}

% ===== Question Block Macro =====
% Usage: \QBlock{Q1}{stem...}{<items>}

\newcommand{\QBlock}[3]{%
  \Needspace{8\baselineskip}%
  \noindent
  \begin{minipage}{\linewidth}
    \QuestionStyle{#1.~#2}

    \vspace{0.25em}
    \begin{enumerate}
      #3
    \end{enumerate}

    \vspace{0.25em}
    \hrule height 0.4pt
  \end{minipage}
  \vspace{0.75em}
}

% ===== Document Start =====

\begin{document}

\begin{center}
{\huge \textbf{Isomorphic Physics Questions}}\\[0.25em]
{\Large Projectile Motion and 2D Kinematics}\\[0.25em]
{\large Multiple choice questions generated below. Answers in RED and incorrect answers in BLUE.}
\end{center}

\vspace{0.75em}

\QBlock{Q1}{A flying squirrel leaps from a tree with initial speed $v_0=12.5\ \mathrm{m/s}$ at an angle of $35^\circ$ away from the vertical (so it is $55^\circ$ above the horizontal). Neglect air resistance. What is the squirrels speed when it reaches its maximum height?}{
  \Wrong{$12.5\ \mathrm{m/s}$}{This assumes the speed stays constant throughout the flight. In ideal projectile motion, only the horizontal component stays constant while the vertical component changes due to gravity.}
  \Wrong{$10.2\ \mathrm{m/s}$}{This comes from using the wrong trigonometric component (e.g., $v_0\cos 35^\circ$), as if the given angle were measured from the horizontal instead of from the vertical.}
  \Correct{$7.2\ \mathrm{m/s}$}{At maximum height the vertical velocity is zero: $v_y=0$. The speed then equals the magnitude of the horizontal component, $v=|v_x|$; since the launch angle is $35^\circ$ from vertical, $v_x=v_0\sin 35^\circ$, so $v=12.5\sin 35^\circ\approx 7.17\ \mathrm{m/s}\approx 7.2\ \mathrm{m/s}$.}
  \Wrong{$0\ \mathrm{m/s}$}{This confuses ``vertical velocity is zero at the top'' with ``total velocity is zero.'' The horizontal component generally remains nonzero, so the squirrel is not at rest.}
}

\QBlock{Q2}{A maglev test cart is moving on a nearly frictionless track with initial velocity $\vec v_0=(28.0\ \mathrm{m/s})\,\hat\imath$. A guidance system then applies a constant acceleration $\vec a=(1.40\ \mathrm{m/s^2})\,\hat\imath+(3.10\ \mathrm{m/s^2})\,\hat\jmath$. Which statement is correct?}{
  \Wrong{The $y$-component of the carts velocity remains constant.}{This would only be true if $a_y=0$. Here $a_y=3.10\ \mathrm{m/s^2}\neq 0$, so $v_y(t)=v_{y0}+a_y t$ changes with time.}
  \Correct{Both the direction of motion and the speed change.}{With constant acceleration, $\vec v(t)=\vec v_0+\vec a\,t$, so $v_x(t)=v_{x0}+a_x t$ and $v_y(t)=v_{y0}+a_y t$. Since both $a_x\neq 0$ and $a_y\neq 0$, both components of velocity change, which changes the direction (the angle of $\vec v$) and also changes the speed $|\vec v|=\sqrt{v_x^2+v_y^2}$ over time.}
  \Wrong{The $x$-component of the carts velocity remains constant.}{This incorrectly assumes $a_x=0$. Because $a_x=1.40\ \mathrm{m/s^2}$, the carts $x$-velocity must change according to $v_x(t)=v_{x0}+a_x t$.}
  \Wrong{None of the above is correct.}{This choice overlooks that option B follows directly from the component equations $v_x(t)=v_{x0}+a_x t$ and $v_y(t)=v_{y0}+a_y t$ when both acceleration components are nonzero.}
}

\end{document}
